% !TEX program = xelatex
\documentclass[11pt,a4paper]{article}

\usepackage[a4paper,margin=2.25cm]{geometry}
\usepackage{fontspec}
\usepackage{polyglossia}
\setmainlanguage{english}
\setmainfont{Times New Roman}
\usepackage{microtype}
\usepackage{amsmath,amssymb,booktabs,tabularx}
\usepackage{caption}
\usepackage{hyperref}
\hypersetup{colorlinks=true,linkcolor=black,urlcolor=black}
\setlength{\parindent}{0pt}
\setlength{\parskip}{0.55em}
\setlength{\emergencystretch}{1.5em}
\renewcommand{\arraystretch}{1.12}

\newcommand{\Hsem}{H_{\mathrm{sem}}}
\newcommand{\AUC}{\mathrm{AUC}}

\title{\vspace{-1.2em}\textbf{Likelihood-weighted Semantic Entropy as a Non-graph Hallucination Baseline:}\\[-0.25em]
Archive-verified Results on RAGTruth QA}
\author{}
\date{}

\begin{document}
\maketitle
\vspace{-1.7em}

\begin{abstract}
This report evaluates likelihood-weighted semantic entropy as a non-graph baseline for response-level hallucination detection on RAGTruth QA.  The method samples 15 new answers from Gemini~2.5~Flash for each original source prompt, clusters the samples by mutual natural-language-inference entailment, and uses entropy of the likelihood-weighted semantic classes as hallucination risk.  A deterministic 560-row source-level manifest was balanced before an established quarantine of source \texttt{12448}; the resulting analysis set contains 559 sources (447 train and 112 held-out test).  All 559 sources were scored and none was unscorable at the fixed 65,535-token output limit.  On the balanced held-out set, the baseline obtained ROC-AUC $0.450$ (bootstrap 95\% interval $[0.345,0.556]$) and F1 $0.667$ ($[0.582,0.742]$).  The train-selected threshold is effectively zero and marks every held-out source as hallucinated (precision $0.500$, recall $1.000$, accuracy $0.500$).  Thus this implementation provides no useful response-level discrimination in the tested direction.  The result is archive-verified: a cache-only replay covered the same 559 sources, produced byte-identical score and unscorable files, and left the cache inventory unchanged.
\end{abstract}

\section{Research Question}

Semantic entropy estimates whether a language model is uncertain about the meaning of answers it would generate for a prompt.  It is an attractive non-graph comparison because it does not construct a knowledge graph or verify answer claims against retrieved evidence.  Following the semantic-uncertainty formulation of Kuhn, Gal, and Farquhar~\cite{kuhn2023}, we test the directional hypothesis that larger semantic entropy of newly sampled answers is associated with a higher probability that the RAGTruth QA source is labelled hallucinated.

The target is deliberately demanding.  RAGTruth labels describe whether one historical response contains an annotated hallucination span~\cite{niu2024}; the semantic-entropy score instead measures the variability of newly generated Gemini answers for the original source prompt.  The baseline should therefore be interpreted as a prompt-level uncertainty signal evaluated against response-level labels, rather than as an evidence-grounded audit of the historical response.

\section{Experimental Protocol}

\subsection{Data, Manifest, and Isolation}

The experiment uses the QA portion of RAGTruth.  The response-level binary label is $y=1$ when the historical response contains at least one annotated hallucination span, and $y=0$ otherwise.  A deterministic, source-level, label-balanced manifest of 560 QA rows was selected with seed 42 under protocol \texttt{ragtruth-qa-source-balanced-v1}.  It contains 448 train and 112 held-out test sources, with 280 sources from each label before quarantine.  Its SHA-256 checksum is
\begin{center}
\small\texttt{4d2239a1cb9713335c0554324f1a337f6bc8790c66d9e06516b04f653d3fba5a}.
\end{center}

Source \texttt{12448} was excluded by the established source-level quarantine policy before analysis.  It belongs to the training partition.  The resulting analysis set therefore contains 447 train sources (223 factual and 224 hallucinated) and 112 held-out test sources (56 factual and 56 hallucinated), for 559 sources overall.  The held-out labels are not used to construct entropy scores or choose the decision threshold.

The historical answers in this manifest were generated by six distinct RAGTruth source models.  The candidate historical answer itself is not supplied to the entropy scorer.  This avoids label leakage but makes the method a cross-generator, prompt-level baseline.

\subsection{Likelihood-weighted Semantic Entropy}

The implementation follows the public TOHA baseline routines~\cite{toha} while adapting only the generation transport to the authenticated Gemini gateway.  For each original RAGTruth source prompt, Gemini~2.5~Flash generates $N=15$ independent answer samples at temperature $1.0$.  The gateway returns selected-token log probabilities.  For sample $i$, the sample log likelihood is the mean selected-token log probability,
\begin{equation}
  \ell_i=\frac{1}{T_i}\sum_{t=1}^{T_i}\log p(x_{it}\mid x_{i,<t},\mathrm{prompt}),
\end{equation}
which avoids assigning a length advantage to short completions.  Normalised sample masses are
\begin{equation}
  p_i=\frac{\exp(\ell_i)}{\sum_{j=1}^{N}\exp(\ell_j)}.
\end{equation}

The local \texttt{microsoft/deberta-v2-xlarge-mnli} model greedily groups answer samples into semantic classes.  Two samples are considered semantically equivalent when neither NLI direction is contradiction and the two directions are not both neutral.  If $z_i$ is the class assigned to sample $i$, the probability mass of class $k$ is
\begin{equation}
  q_k=\sum_{i:z_i=k}p_i,
\end{equation}
and the reported score is
\begin{equation}
  \Hsem=-\sum_k q_k\log q_k.
\end{equation}
Larger $\Hsem$ is pre-specified to mean greater hallucination risk.  The sampling cap is fixed at Gemini~2.5~Flash's documented maximum of 65,535 output tokens.  An answer reaching that cap would be explicitly recorded as an output-length-unscorable outcome, rather than silently dropped or treated as a low-risk score.

\subsection{Evaluation and Reproducibility Controls}

The only learned decision parameter is a threshold $\theta$, selected on the 447-source training partition by maximum F1 under the rule $\Hsem\geq\theta\Rightarrow\widehat{y}=1$.  The held-out partition is evaluated once after this choice is frozen.  ROC-AUC is threshold-free; F1, precision, and recall use the train-selected threshold.  The reported 95\% intervals are percentile intervals from 1,000 non-parametric bootstrap resamples of the held-out set (seed 42 for ROC-AUC and seed 43 for F1).

The logical generation model is \texttt{openai/gemini-2.5-flash}; the archive pins the authenticated gateway-manifest SHA-256, runtime fingerprint, and selected-token-logprob protocol.  Raw prompts, completions, and credentials are absent from the report and terminal archive.  Content-addressed external-disk caches retain the material needed for deterministic cache-only replay.

\section{Results}

\subsection{Coverage and Held-out Metrics}

All eligible sources completed: 559/559 were scored and 0/559 were marked output-length-unscorable.  Table~\ref{tab:metrics} reports the single held-out evaluation.  The F1 threshold selected on train is $\theta=-1.0\times10^{-9}$, numerically just below zero; its training F1 is $0.6677$.

\begin{table}[h]
\centering
\caption{Held-out results on 112 balanced test sources (56 factual and 56 hallucinated).  Intervals are 1,000-resample percentile bootstrap 95\% intervals.}
\label{tab:metrics}
\small
\begin{tabular}{@{}lc@{}}
\toprule
Quantity & Value \\
\midrule
ROC-AUC & $0.4503\;[0.3454,0.5555]$ \\
Precision & $0.5000$ \\
Recall & $1.0000$ \\
F1 & $0.6667\;[0.5823,0.7416]$ \\
Accuracy & $0.5000$ \\
TP / FP / FN / TN & $56$ / $56$ / $0$ / $0$ \\
\bottomrule
\end{tabular}
\end{table}

The threshold predicts the hallucination class for all 112 held-out sources.  Consequently, the reported F1 is exactly the value achieved by an all-positive classifier on a balanced set; it is not evidence of useful discrimination.  The ROC-AUC is below $0.5$, and its bootstrap interval contains $0.5$.

Table~\ref{tab:distribution} gives score and semantic-class summaries without exposing response text or generated completions.  In this held-out sample, factual sources have a \emph{higher} mean and median entropy than hallucinated sources.  This is the direction opposite to the pre-specified interpretation of higher entropy as greater hallucination risk, consistent with the observed ROC-AUC below chance.

\begin{table}[h]
\centering
\caption{Held-out descriptive score statistics.  ``Classes'' is the number of mutual-entailment semantic classes among 15 Gemini samples.}
\label{tab:distribution}
\small
\begin{tabular}{@{}lccccc@{}}
\toprule
Label & $n$ & Mean $\Hsem$ & Median $\Hsem$ & Mean classes & Median classes \\
\midrule
Factual ($y=0$) & 56 & 0.387 & 0.320 & 1.98 & 2 \\
Hallucinated ($y=1$) & 56 & 0.311 & 0.252 & 1.86 & 2 \\
\bottomrule
\end{tabular}
\end{table}

\subsection{Archive Acceptance}

The terminal archive verifies the full live stage and a separate cache-only acceptance stage.  The latter completed all 559 sources with zero live inference calls.  The live and replay \texttt{scores.jsonl} files are byte-identical; the corresponding \texttt{unscorable.jsonl} files are also byte-identical.  Finally, the pre- and post-replay cache inventories are byte-identical.  These checks establish reproducibility of this fixed run; they do not turn the run into an independent second evaluation.

\section{Interpretation}

The tested hypothesis is not supported on this manifest.  In its prescribed direction, likelihood-weighted semantic entropy has no useful ability to separate hallucinated from factual historical RAGTruth responses.  The selected threshold collapses to the all-positive rule, yielding perfect recall but a false-positive rate of one.  Its AUC of $0.450$ also indicates that, as a ranking signal, larger prompt-level entropy tends to occur more often for the factual class in this sample.

This outcome is informative rather than a failure of experimental control.  Semantic entropy measures dispersion among answers that \emph{Gemini} would produce from a prompt.  RAGTruth's label instead describes the factual support of one historical answer, generated by a different source model, relative to its supplied context.  A prompt can admit several semantically different yet context-supported answers, thereby increasing entropy without making the historical answer hallucinated.  Conversely, a model can repeatedly produce the same unsupported answer, producing low semantic entropy while the historical response is hallucinated.  The observed reversal is consistent with this construct mismatch, but the present experiment cannot establish its causal mechanism.

The result should therefore be used as a negative non-graph baseline under the stated protocol, not as evidence that semantic entropy is generally ineffective.  It also should not be numerically compared as a method improvement or degradation against the graph reports: those reports use a different 750-row manifest, different scorable denominators, and response-conditioned evidence signals.  A valid paired comparison would require a common manifest and a pre-specified evaluation protocol for both method families.

\section{Limitations and Next Tests}

This is one deterministic manifest and one sampling configuration ($N=15$, temperature $1.0$) evaluated with one generator.  The bootstrap intervals quantify sampling variability of the fixed held-out set; they are not a paired significance test and do not establish generalisation to other prompts, models, or sampling temperatures.  The local NLI equivalence rule is itself a modelling choice, and semantic-class assignments may be imperfect.

Most importantly, this baseline is intentionally answer-agnostic: it never reads the response whose label is being predicted.  It cannot distinguish a factually correct response from a hallucinated response that arose under the same source prompt except through correlations between prompt-level uncertainty and labels.  A future uncertainty baseline that conditions on the candidate answer, or a pre-registered paired evaluation on the 750-row manifest, would test a different and more response-aligned hypothesis.  Such an extension should retain the same train-only threshold rule and cache-only acceptance criterion.

\section{Conclusion}

On the archive-verified 559-source RAGTruth QA evaluation, likelihood-weighted semantic entropy is a completed, reproducible, non-graph baseline but not a useful detector of the historical response-level hallucination label.  Complete coverage and cache-only replay establish the integrity of the result.  The held-out ROC-AUC of $0.450$ and all-positive train-selected decision rule reject the baseline's intended higher-entropy/higher-risk association for this protocol.  The result is best understood as evidence that prompt-level generation uncertainty is not an adequate substitute for response-grounded hallucination assessment in this setting.

\clearpage
\begin{thebibliography}{9}

\bibitem{kuhn2023}
Lorenz Kuhn, Yarin Gal, and Sebastian Farquhar.
\newblock Semantic Uncertainty: Linguistic Invariances for Uncertainty Estimation in Natural Language Generation.
\newblock arXiv:2302.09664, 2023.

\bibitem{niu2024}
Cheng Niu, Yuanhao Wu, Juno Zhu, Siliang Xu, KaShun Shum, Randy Zhong, Juntong Song, and Tong Zhang.
\newblock RAGTruth: A Hallucination Corpus for Developing Trustworthy Retrieval-Augmented Language Models.
\newblock In \emph{Proceedings of the 62nd Annual Meeting of the Association for Computational Linguistics (Volume 1: Long Papers)}, pages 10862--10878, 2024.
\newblock \url{https://doi.org/10.18653/v1/2024.acl-long.585}.

\bibitem{toha}
SB AI Lab.
\newblock TOHA: Hallucination Detection in LLMs with Topological Divergence on Attention Graphs (public reference implementation).
\newblock \url{https://github.com/sb-ai-lab/TOHA}, accessed August 2026.

\end{thebibliography}

\end{document}
